\section{Two-stage evaluation and implementation}\label{app:implementation}

\subsection{Two-stage evaluation}\label{sec:twostage}

Most of the integration budget goes to trajectories that settle early onto a known regime yet are integrated over the full horizon, which the protocol requires only for target candidates (section~\ref{sec:protocol}). We therefore evaluate every initial condition in two stages, for all strategies, including uniform sampling. Stage~1 integrates the transient segment up to the end of the early window, i.e., for $\min(n_{\rm tr},n_{\rm obs})$ steps, and then a short observation window of $n_{\rm obs}/4$ steps; here $n_{\rm tr}$ and $n_{\rm obs}$ are the numbers of steps of the transient segment and of the observation window. The trajectory stops after stage~1 if, on this window, it is identified with confidence as a known nontarget regime: an equilibrium (converged, or decaying with an amplitude below $0.02L$), or a known cycle or chaotic attractor with a $z$ score below half the identification tolerance. Escaping trajectories are final in any case. All other trajectories (target candidates and unknown or ambiguous ones) continue from the end of the early window to the full horizon. Since stage~2 resumes the same trajectory at the end of the early window, the fingerprints of continued trajectories coincide bit for bit with those of a single full integration, and the target is always confirmed by the full protocol. A stopped trajectory keeps its label; only its $\xi$ is computed on the shorter window.

For continued trajectories, stage~1 costs extra steps, because a quarter of the observation window is integrated twice. A rule based on step counts therefore keeps the two-stage mode only while it pays off. After the first 200 evaluations, the mode stays on as long as the steps saved over the last ten batches exceed the extra steps; otherwise, only every tenth batch is evaluated in two stages, as a probe. The rule uses step counts rather than time, so runs remain reproducible. On uniform samples from eight problems, the two-stage labels and the fingerprints of all continued and escaping trajectories coincided with those of full integration bit for bit, and no target hit was lost (section~\ref{sec:cost}).

\subsection{Implementation details}\label{sec:impl}

The computations of the arms themselves (nearest-neighbor queries of the active-learning arm, the chains, and the matching) take 1--8\% of the wall time on the Chua systems with up to two copies, 28\% on three uncoupled copies of Chua~II, and 20--60\% on the cheap problems, where an integration costs microseconds. For $n\ge6$, brute-force neighbor search speeds up the queries by a factor of 1.6--2.6 relative to a tree ($G_T$ rises from 16 to 29 on three coupled copies of SE11). The cap on the training sample of the classifier (the last $2\times10^4$ evaluated points) removes the growth of the overhead in long runs. On three uncoupled copies of Chua~II, it lowers the overhead from 51\% to 28\%; by training the classifier on recent points only, it also lowers the median cost from $1.4\times10^5$ to $8.6\times10^4$ integrations and raises the number of runs that succeed within $4\times10^5$ integrations from 8 to 9 out of 10. A wider batch, with two or four times more chains and boundary points to occupy more cores, does not pay off: the number of evaluations grows by a factor of 1.5--3, whereas the time per evaluation decreases by only about 10\%. The time gains of Table~\ref{tab:cost} were measured with four threads per process, and those of Table~\ref{tab:time} on one thread. With the batch of 500 evaluations used here, the portfolio parallelizes better than uniform sampling on the coupled Chua circuits: its time to detection falls 4.0-fold from one to four threads, against 1.6-fold for the cost per uniform evaluation. There, the four-thread gain is therefore larger (68 against 27), and the one-thread value is the conservative one. On Chua~II, whose runs are short, the two values are close (4.1 and 5.7).

\subsection{Comparison with recurrence-based stopping}\label{sec:recurrence}
\begin{table}[tbp]
\centering
\caption{Uniform sampling with three evaluation schemes against the portfolio, on one thread ($R=20$): cost of one evaluation (ms), cost $N_{0.9}$ and time to detection $T_{0.9}$ (s) of the portfolio, and its time gain relative to uniform sampling with two-stage evaluation ($G_T$) and with recurrence-based stopping ($G_T^{\rm rec}$), with 95\% bootstrap intervals. For uniform sampling, $T_{0.9}$ is the fixed reference cost $N^0_{0.9}$ of section~\ref{sec:measure} times the cost of one evaluation; for the portfolio, it is the measured 0.9 quantile of the run time. Attractors.jl: the fastest of the tested settings. Bold: values whose lower interval bound exceeds 1.}
\label{tab:time}
\footnotesize
\setlength{\tabcolsep}{3pt}

\begin{tabular}{lrrrrrrrr}
\toprule
Problem & $N^0_{0.9}$ & \multicolumn{3}{c}{Cost per evaluation, ms} & \multicolumn{2}{c}{Portfolio} & $G_T$ & $G_T^{\rm rec}$ \\
 & & full & two-stage & \shortstack{Attr-\\actors.jl} & $N_{0.9}$ & $T_{0.9}$, s & & \\
\midrule
%TABTIME-BEGIN
Chua II & 818 & 26.7 & 6.07 & 1.95 & 113 & 0.87 & \textbf{5.7} & \textbf{1.8} \\
 &  &  &  &  &  &  & \ci{4.7}{9.1} & \ci{1.5}{2.9} \\[1pt]
2 Chua II, $\epsilon=0.1$ & $1.05\times10^{5}$ & 72.8 & 44.8 & 16.4 & 1899 & 172 & \textbf{27.3} & \textbf{10.0} \\
 &  &  &  &  &  &  & \ci{11}{48} & \ci{4.2}{18} \\[1pt]
SE11 & 317 & 0.131 & 0.0562 & 0.131 & 71 & 0.028 & 0.64 & 1.5 \\
 &  &  &  &  &  &  & \ci{0.4}{0.71} & \ci{0.93}{1.6} \\[1pt]
LE1 & 1617 & 0.021 & 0.0108 & 0.0819 & 413 & 0.58 & 0.03 & 0.23 \\
 &  &  &  &  &  &  & \ci{0.0057}{0.17} & \ci{0.043}{1.3} \\[1pt]
%TABTIME-END
\bottomrule
\end{tabular}

\end{table}

The cost of one evaluation can be reduced further by stopping a trajectory as soon as it recurs in the cells of an already known attractor, as in Attractors.jl~\cite{DatserisWagemakers2022,Datseris2023Framework}. For uniform sampling, the index of the first point in the target basin does not depend on how the trajectories are evaluated. Hence the time of uniform sampling with any evaluation scheme is $N^0_{0.9}$ times the mean cost of one evaluation, and the schemes can be compared on a common set of initial conditions. Table~\ref{tab:time} compares, on one thread, full integration, our two-stage evaluation, and \textsc{AttractorsViaRecurrences} (Attractors.jl 2.0; a sparse grid of 100 cells per dimension, a sampling interval of 0.1, and adaptive Tsit5 with tolerance $10^{-9}$ for the Chua systems and fixed-step RK4 with our step for SE11 and LE1, whichever was faster). The labels of the recurrence mapper on the first 500 points coincide with ours, and it recognizes the target states fed to it. On the bounded Chua systems with long transients, recurrence-based stopping is about 3 times cheaper per evaluation than our two-stage evaluation, which still needs the early window and a quarter of the observation window for the fingerprint. On the escaping systems, our evaluation is 2--8 times cheaper, because the escape radius detects escape at once. Against uniform sampling with recurrence-based stopping, the time gain of the portfolio falls from 27 to 10 on the coupled Chua circuits and from 5.7 to 1.8 on Chua~II, both still significant; the gain in integrations is unaffected. The search strategy and the evaluation back end are independent choices: the portfolio needs only a regime label and a transient duration for each point, both of which the recurrence mapper provides (section~\ref{sec:discussion}).

